%!TEX root = ../main.tex
% =============================================================================
%  第 4 章  模型假设
%  本节含：4.1 假设表 A1--A12（编号/内容/依据/验证方式）｜4.2 信息结构声明 R1/R2/R3
%    4.1 假设表（A1--A12：编号／假设内容／依据／验证方式；"依据"列优先引用第 3 章的核查结论）
%    4.2 信息结构声明（R1/R2/R3 三种解读与本文主方案）
%  判据：一条假设若不能回答"不用它模型会怎样"，它就不该出现在假设表里；
%        每条假设的"验证方式"必须能在第 6--9 章的 x.3 里兑现。
% =============================================================================

\section{模型假设}\label{chap:assump}

\subsection{假设表}

题面已给定储能最大容量 \(12\,000\)~kWh、最大充放电功率 \(5000\)~kW、2025 年 1 月 1 日
\(0{:}00\) 储电量 \(6000\)~kWh、储电量运行区间 \(1200\)--\(10800\)~kWh 与充放电效率 \(90\)\%；
本文取运行区间为额定容量的 \(10\)\%--\(90\)\%（可用摆幅 \(9600\)~kWh），二者与功率上限构成两类硬
约束，谁更紧由\chapref{chap:q1}的容量/功率双向扫描判定。除此之外本文引入 A1--A12 十二条假设，逐条的\textbf{必要性、
现实或题目依据、影响方向与验证方式}见表~\ref{tab:assumptions}；\textbf{判据是一条假设必须能回答
“不用它模型会怎样”}，装饰性假设一律删除。

\begin{table}[htbp]
  \centering
  \caption{模型假设及其必要性、现实/题目依据、影响方向与验证方式（A1--A12）}
  \label{tab:assumptions}
  \scriptsize
  \setlength{\tabcolsep}{3pt}
  \begin{tabularx}{\linewidth}{C{0.62cm} L{2.65cm} L{4.30cm} L{1.95cm} Y}
    \toprule
    编号 & 假设 & 必要性（不用会怎样）与依据 & 影响方向 & 验证或灵敏度方式 \\
    \midrule
    A1 & 不允许向外部电网售电
      & 否则目标多出售电收益项、盈余有了出路、弃光恒为 0，可行域与最优值都改变；依据：题面未给上网电价，
        政策导向“优先自用”
      & 费用偏高，结果保守
      & 非负性审计（审计项 AU7）五个结果文件负购电量数均为 0；另对上网电价 \(\theta p_t\) 作政策情景扫描 \\
    A2 & 充电、放电效率\textbf{单边各取} \(0.9\)（往返 \(0.81\)）
      & 效率决定套利是否有利可图，是收益的核心参数；依据：题面只说“充放电效率为 \(90\%\)”，
        \textbf{未区分充、放方向}
      & 效率越低，往返损耗越大，收益越低
      & 四种读法对照（单边各 \(0.9\)／往返 \(0.9\) 等效单边／仅充电侧／仅放电侧），节省额极差
        \(13.23\%\)（表~\ref{tab:sens-q1}） \\
    A3 & 同一时段不同时充放电
      & 否则须引入 \(0\)--\(1\) 变量，模型由 LP 变为 MIP，规模上升且对偶信息不再可用；依据：
        \(\eta_c\eta_d<1\) 时同时充放严格劣于只做净动作（命题~\ref{prop:no-simul}）
      & 不影响最优值，只影响求解方式
      & 断言 \(\min(c_t,d_t)=0\) 全部成立：问题 1--4 的互斥违反数均为 0（审计项 AU5） \\
    A4 & 忽略自放电与循环衰减
      & 否则目标须加入与放电量相关的衰减成本项并需衰减参数；依据：题面未给任何衰减参数、窗口仅 334 天
      & 高估长期收益，放电量越大高估越多
      & 加入循环成本 \(c_{cyc}\)（元/kWh 放电量）作单参数灵敏度，给出“储能节省额归零”的临界值 \\
    A5 & 储电量\textbf{跨日连续递推}，不逐日清零
      & 逐日清零与“2025-01-01 \(0{:}00\) 储电量 \(6000\)~kWh”的初值条件矛盾，并使每日初值成为
        人为参数；依据：题面只要求 \(0{:}00\) 与 \(24{:}00\) 储电量相同。
        \textbf{诚实说明}：由于 A6 的每日终端条件 \(S_{144}\ge S_0\) 与 \(S_0\) 逐日不变，实测 334 天的
        \(24{:}00\) 储电量\textbf{每日恰等于 \(6000.00\)~kWh}（跨日带电量 \(0\)~kWh），故本窗口内
        \emph{并未}发生跨日套利；A5 与“每日重置初值”在本题数据上的差别因此可忽略，
        它的实际作用是使初值有来源、口径可复现，而不是释放跨日收益
      & 影响全窗口费用的记账方式；本窗口实测与“每日重置初值”等价
      & 与“每日重置初值”变体对照；审计项 AU6 逐日核对“末值不低于初值”的违反天数（问题 2--4 均为 0）
        与跨日带电量（实测恒为 \(0\)~kWh） \\
    A6 & 每日终端条件 \(S_{144}\ge S_0\)；问题 1 取严格日周期 \(S_0=S_{144}=6000\)~kWh
      & 否则模型会在每天末尾与窗口末尾“掏空”储能压低当期费用，形成不可执行的短视解；依据：储能运营
        常识与题面“两端储电量相同”
      & 结果保守；日周期比自由环状贵 \(8.35\) 元（\(0.02\%\)）
      & 终端处理四组对照：完全自由、\(S_T\ge S_0\)、日周期、水价值终端（表~\ref{tab:sens-q1}） \\
    A7 & 调整决策“后覆盖前”，计价基准恒为 \(0{:}00\) 的\textbf{计划购电量}
      & 否则调整费与违约费口径不可复现，也无法与题面表述对应；依据：题面“计划购电量高于调整购电量的
        部分”一句明确了基准
      & 影响调整与违约相关费用
      & 与“相对上一次调整量”口径并列报告并各自命名（\calref{C-09}） \\
    A8 & 2025 年 1 月作为储电量递推的\textbf{预热期}，其终值即 2 月 1 日 \(0{:}00\) 初值
      & 结果窗口自 2025-02-01 起而初值自 2025-01-01 起算，不作声明则 2 月 1 日储量没有来源；依据：
        附件 2 覆盖全年、结果文件要求 334 天
      & 影响窗口初期费用；1 月不计入结果
      & 论文与结果文件中明示 1 月不计入并公开初值来源；审计项 AU6 逐日核对终端条件 \\
    A9 & \textbf{信息结构取 R3}：问题 2/4-2 的计划依据为附件 3 当日 \(0{:}00\) 预报，负载按“已知”
      处理；问题 3/4-3 另有三次滚动的当次预报
      & 题面未说明 \(0{:}00\) 知道什么，不声明则问题 2--4 的任何数值都无法解释（三种解读的窗口总费用
        相差近千万元）；依据：题面“根据预测数据制定计划”与附件 3 的发布时刻
      & 费用落在 R1（下界）与 R2（上界）之间；负载按已知处理\textbf{低估}紧急购电量，结果偏乐观
      & R1/R2/R3 在同一 334 天窗口与同一场景集上的完整对照（表~\ref{tab:infostruct}）；
        \textbf{非预期性置换检验}：置换任一次预报都不改变早于其发布时刻的决策（\S9.2） \\
    A10 & 波动电价下 \(0{:}00\) 可用的\textbf{电价信息取“次日可预知的分时骨架”}；主口径的首阶段电价
         系数用\textbf{情景概率加权均值} \(\bar p^{\text{sc}}_t\)（代表日实测电价的层加权平均，
          两类信息都不含当日实际电价）；结算仍按附件 4 实际电价
      & 附件 4 未给日前电价预报；不声明则“计划”会用到结算时的实际电价，构成数据泄漏。\textbf{主口径实际
        用的是 \(\bar p^{\text{sc}}_t\)}，与分时骨架的最大偏差 \(0.0462\)~元/kWh（占均价 \(6.02\%\)）；
        本文把两者作为“骨架／历史加权均值哪个更适合做日前代理”的并列口径报告，而非只声明其中一个；
        依据：骨架与附件 4 逐时段均值最大偏差仅 \(5.15\times10^{-5}\)~元/kWh
      & 忽略电价残差波动，使问题 4 结果\textbf{偏乐观}；两种代理的偏差已量化
      & 结算侧改用实际电价；\textbf{调整层同样用 \(\bar p^{\text{sc}}_t\)（不是当日实际价，故无数据泄漏）}；
        并以“是否保留残差自相关”作对照（分块 Bootstrap 的分位数最大相对偏差达 \(20.44\%\)） \\
    \midrule
    A11 & \textbf{“每天的电价相同”理解为附件 1 的 144 点日内电价曲线逐日重复}：日内电价仍随时间
         变化，仅排除电价在日与日之间的波动
      & 附件 1 的 144 个电价中 \(137\) 个互不相同、区间 \(0.3713\)--\(1.3952\)~元/kWh（峰谷比
        \(3.76\)），故它是\textbf{日内曲线}而非单一常数；题面背景段明确“在电价较低的时间段充电、
        并在电价较高的时间段释放电能”，若日内恒价则该机制无对象可实现
      & 若误读为\textbf{全天单一电价}，日内套利项消失，问题 2 的储能价值只剩消纳提升，节省率随之改变
      & 主结果按日内曲线计算；“全天单一均价”仅作\textbf{反面对照}并给出代价（\S9.3） \\
    \midrule
    A12 & \textbf{充放电量在 \(0{:}00\) 一次性确定、实时不允许再调度}；计划量买到的电不可转手，
          超购部分既不产生收益也不再被后续时段使用（\textbf{超购即弃}）
      & \textbf{本文影响输出量级最大的结构性假设}：题面只锁定\emph{结算}口径（“其他时间段的购电费用
        均按计划购电量计算”）与“\(0{:}00\) 制定计划”，并未禁止储能按实际值动作。若允许实时再调度，
        缺口可由电池直接补足、盈余可被电池吸纳，\(5p_t\) 紧急购电与盈余处置量都会大幅下降；
        不声明则该项代价无法与模型结论分离。依据：题面把“缺口 \(\to\) 5 倍紧急购电”设为主机制，
        隐含“计划先定、缺口事后补”的读法，故本假设是可辩护的工程选择而非题面强制
      & 结论\textbf{偏保守}：费用偏高、紧急购电与弃电偏多；不改变“储能降低购电费”的方向
      & \textbf{实时再调度消融（表~\ref{tab:abl-q2} 的 RT 行）}：固定已结算的 \(E^{\text{plan}}\)，
        用当日实际净需求重优化 \(c,d,S\)，量化紧急购电与盈余处置量的差额 \\
    \bottomrule
  \end{tabularx}
\end{table}

A1--A12 十二条中，\textbf{A12 决定两项标志性输出的量级}（紧急购电与盈余处置量），RT 消融给出它的
完整代价；\textbf{A2 影响节省率最大}，其四种读法的量化对照见式~\eqref{eq:q1-eff} 与
灵敏度表~\ref{tab:sens-q1}；主口径写入全部结果表，其余读法只在灵敏度节出现（\evid{C-Q1-03}）。

\subsection{信息结构声明}

问题 2--4 中“每天 \(0{:}00\) 制定计划”有三种解读。本文各跑一次全窗口实验，条件完全一致（同一
334 天窗口、同一场景集、同一终端条件、同一电价口径），只有“\(0{:}00\) 已知什么”不同，
三行可直接相减（表~\ref{tab:infostruct}）。本文采用\textbf{信息结构 R3}。

\begin{table}[htbp]
  \centering
  \caption{信息结构的三种解读及其在 2025-02-01 至 12-31（334 天）上的实测后果
           （费用单位：元，电量单位：kWh，电价均为附件 1 固定价）}
  \label{tab:infostruct}
  \footnotesize
  \setlength{\tabcolsep}{3pt}
  \begin{tabularx}{\linewidth}{C{0.7cm} L{2.95cm} C{2.25cm} C{1.35cm} C{2.15cm} Y}
    \toprule
    代号 & 解读 & 总费用 & 相对基线节省率 & 紧急购电量 & 本文的使用方式 \\
    \midrule
    R1 & 完全信息：\(0{:}00\) 即知当日负载与光伏真值
      & \(13\,411\,560.90\) & \(18.26\%\) & \(0.00\)（0/334 天）
      & 不作主方案；用作“预测误差的代价”的\textbf{下界对照} \\
    R2 & 典型日预报：以附件 1 的日均形状作为预报
      & \(23\,484\,646.63\) & \(-43.14\%\) & \(2\,525\,548.76\)（286/334 天）
      & 仅作悲观对照：附件 1 是实际值的全年均值而\textbf{不是预报}，本文\textbf{不采用} \\
    R3 & 日前预报：附件 3 当日 \(0{:}00\) 发布的光伏预报，负载按已知处理
      & \(14\,264\,792.13\) & \(13.06\%\) & \(172\,015.82\)（296/334 天）
      & \textbf{本文主方案}：最符合工程实际与题面“制定计划”的语义 \\
    \bottomrule
  \end{tabularx}
\end{table}

R3 的三点补充：① 计划依据（附件 3 预报 \(+\) 负载曲线）与结算依据（附件 2 当日实际负载与光伏）
是\textbf{两套数据}，差额即紧急购电来源；② “预测误差的代价”为
\(14\,264\,792.13-13\,411\,560.90=853\,231.23\) 元、占 R1 费用的 \(6.36\%\)（\evid{C-Q2-01}）；
③ 负载按“已知”处理是对 R3 的乐观简化（题面未给负载预报），会低估紧急购电量，
\appref{sec:verify-three-app}{D.1 三层验证与审计结论}另作预报误差灵敏度。
